\documentclass[12pt]{article}
\usepackage[utf8]{inputenc}
\usepackage[T1]{fontenc}
\usepackage{hyperref}
\title{Isubu Language Resources for K-12 Curriculum Development}
\author{Compiled by Claude Sonnet 5.5 (claude-sonnet-5-5) for Andr\'as Kornai}
\date{}
\begin{document}
\maketitle

\section{Introduction}
Isubu (ISO 639-3 szv, also called Su, Isu, Isuwu, Subu or Bimbia) is a Bantu language (Duala group, A.23) spoken in the Limbe and Tiko area of Fako Division, South West Region, Cameroon. It is endangered and is being replaced by Cameroon Pidgin English. Estimates of speakers vary, and one source gives roughly 600. It should not be confused with Isu, a Grassfields language. This report is based on a single round of web searches, which found very little.

\section{Primary Educational Resources}
No primer, textbook or school programme was found.

\section{Linguistic and Reference Materials}
The DOBES documentation project covers Bubia and Isubu, including phonology, morphology, syntax and pragmatics, because neither language has an adequate description. See \url{https://dobes.mpi.nl/?p=1024}. The endangered languages catalogue entry is at \url{https://endangeredlanguages.com/node/91997}. It lists no standard writing system. A 2020 book on plant and animal vocabulary in Isubu and Wovea was mentioned in the search summary, but I saw no URL for it. No dictionary was found.

\section{Religious and Community Resources}
General background on the Subu people is at \url{https://en.wikipedia.org/wiki/Subu_people}. No scripture or community organisation was found.

\section{Digital and Audio Resources}
The DOBES project presumably holds recordings, but the search results did not confirm their content.

\section{Curriculum Development Context}
The language has no known standard orthography and no known school use. A curriculum would first need community orthography work, based on the DOBES documentation.
\end{document}